%% 50장 분포족
\chapter{분포족}
\label{ch:distribution-families}

\abstract{통계 라이브러리가 다루는 분포는 서로 떨어진 목록이 아니라 몇 가지 구성으로 이어진 하나의 가족이다. 정규 확률변수의 제곱합이 카이제곱이 되고 그 비가 $t$와 $F$가 되며 감마 확률변수의 비가 베타가 된다. 이항은 Poisson으로, Poisson과 감마의 혼합은 음이항으로 이어진다. 이 장에서는 stardust가 구현한 단변량 분포 서른한 개를 이런 관계로 엮어 소개한다. 확률변수의 변환이 밀도를 어떻게 바꾸는지부터 시작해, 정규에서 나오는 표본 분포들의 밀도를 유도하고 감마와 베타의 관계, 이항의 누적분포함수가 불완전 베타 함수이고 Poisson의 누적분포함수가 불완전 감마 함수라는 항등식을 증명한다. 이 항등식들이 40장의 특수 함수가 통계의 계산 엔진인 까닭이다. 이어서 혼합으로 만드는 분포, 최댓값의 극한에서 나오는 극값 분포, 검정력 계산에 쓰이는 비중심 분포를 다루고 비중심 카이제곱의 Poisson 혼합 급수에서 실제로 몇 항이 필요한지를 확인한다.}

\section{변환과 위치--척도족}
\label{sec:transformations}

확률변수 $X$의 밀도가 $f_X$이고 $g$가 열린 구간에서 미분 가능하고 $g'>0$이면 $Y=g(X)$의 누적분포함수는 $F_Y(y)=P(X\le g^{-1}(y))=F_X(g^{-1}(y))$이고 미분하면
\[
f_Y(y)=f_X\bigl(g^{-1}(y)\bigr)\,\bigl|(g^{-1})'(y)\bigr|
\]
이다. 순감소하는 $g$도 절댓값 덕분에 같은 식이다. 다변수에서는 미분의 절댓값이 야코비안 행렬식의 절댓값으로 바뀐다. 이 식의 가장 단순한 경우가 $Y=\mu+\sigma Z$($\sigma>0$)로, $f_Y(y)=\sigma^{-1}f_Z((y-\mu)/\sigma)$이다. 하나의 표준 분포 $Z$에서 이렇게 얻는 분포들을 \term{위치--척도족}{location-scale family}이라 한다. 정규, 로지스틱, 라플라스, 코시, Gumbel, 균등이 모두 위치--척도족이고 계산은 표준화한 $z=(x-\mu)/\sigma$에서 한다. 39장의 $e^{-x^2}$에서 본 것처럼 지수 앞의 반올림이 크게 증폭되므로 이 표준화 자체가 오차의 원천이 될 수 있어서 stardust는 $z$를 double-double로 만든다(\cref{sec:families-numerics}).

음이 아닌 지지집합의 분포에서는 척도만 있는 \term{척도족}{scale family} $Y=sX$를 쓴다. 지수, 감마, 베이불, 레일리, 파레토, 역감마가 그렇고 stardust는 이들의 척도 매개변수를 모두 \code{scale}로 부른다. 로그정규는 $\log X$가 정규인 분포로, 위치--척도족을 $\exp$로 옮긴 것이다.

\section{정규에서 나오는 분포}
\label{sec:normal-derived}

\subsection{카이제곱}

$Z_1,\dots,Z_k$가 독립인 표준 정규이면 $Q=\sum Z_i^2$의 분포를 자유도 $k$의 \term{카이제곱 분포}{chi-square distribution}라 한다. 밀도를 구하자. $k=1$이면 $P(Z^2\le q)=2\Phi(\sqrt q)-1$을 미분해 $f(q)=q^{-1/2}e^{-q/2}/\sqrt{2\pi}$이다. 이것은 모양 $1/2$, 척도 $2$인 감마 분포의 밀도
\[
f(x;a,s)=\frac{x^{a-1}e^{-x/s}}{\Gamma(a)\,s^a},\qquad x>0
\]
이다($\Gamma(1/2)=\sqrt\pi$). 감마 분포는 척도가 같으면 모양이 더해진다. 독립인 $X\sim\mathrm{Gamma}(a,s)$, $Y\sim\mathrm{Gamma}(b,s)$의 합의 밀도는 합성곱
\[
\int_0^z\frac{x^{a-1}(z-x)^{b-1}e^{-z/s}}{\Gamma(a)\Gamma(b)s^{a+b}}\,dx=\frac{z^{a+b-1}e^{-z/s}}{\Gamma(a)\Gamma(b)s^{a+b}}B(a,b)=f(z;a+b,s)
\]
이기 때문이다($x=zt$로 치환하면 베타 적분 $B(a,b)=\Gamma(a)\Gamma(b)/\Gamma(a+b)$가 나온다). 그래서 $\chi^2_k=\mathrm{Gamma}(k/2,2)$이고 카이제곱의 누적분포함수는 40장의 정규화된 불완전 감마 함수 $P(k/2,x/2)$이다. 지수 분포는 모양 $1$인 감마이므로 $\chi^2_2$는 평균 $2$인 지수 분포이다.

\subsection{$t$와 $F$}

$Z$가 표준 정규, $V$가 $\chi^2_\nu$이고 둘이 독립이면
\[
T=\frac Z{\sqrt{V/\nu}}
\]
의 분포를 자유도 $\nu$의 \term{$t$ 분포}{Student's t distribution}라 한다. 정규 모집단에서 표본평균과 모평균의 차를 표준오차로 나눈 값의 분포는 Gosset(필명 Student)이 1908년에 처음 구했고\cite{student1908}, 지금의 $t$ 꼴과 엄밀한 증명은 Fisher가 정리했다. $V=v$로 조건을 걸면 $T$는 분산 $\nu/v$인 정규이므로 $T$의 밀도는 정규 밀도를 $V$의 분포로 적분한 것이고 그 적분은 감마 적분이 되어
\[
f(t;\nu)=\frac{\Gamma\bigl((\nu+1)/2\bigr)}{\sqrt{\nu\pi}\,\Gamma(\nu/2)}\Bigl(1+\frac{t^2}\nu\Bigr)^{-(\nu+1)/2}
\]
이다. 꼬리가 $|t|^{-(\nu+1)}$로 줄어 $\nu$차 이상의 모멘트가 없다. $\nu=1$이면 코시 분포이다. $\nu\to\infty$이면 큰 수의 법칙으로 $V/\nu\to1$이어서 $T$는 표준 정규로 수렴하고 밀도도 $(1+t^2/\nu)^{-(\nu+1)/2}\to e^{-t^2/2}$로 수렴한다.

$U\sim\chi^2_m$, $V\sim\chi^2_n$이 독립이면 $(U/m)/(V/n)$의 분포를 \term{$F$ 분포}{F distribution}라 한다. 분산분석에서 귀무가설 아래 두 분산 추정값의 비가 이 분포를 따른다. $T^2=Z^2/(V/\nu)$이므로 $T^2\sim F(1,\nu)$이고 $n\to\infty$이면 $V/n\to1$이어서 $mF(m,n)$은 $\chi^2_m$으로 수렴한다. $F$의 누적분포함수는 다음 절의 베타 관계로 불완전 베타 함수가 된다.

\section{감마와 베타}
\label{sec:gamma-beta}

\begin{proposition}
\label{prop:gamma-beta}
$X\sim\mathrm{Gamma}(a,s)$, $Y\sim\mathrm{Gamma}(b,s)$가 독립이면 $B=X/(X+Y)$는 \term{베타 분포}{beta distribution} $\mathrm{Beta}(a,b)$, 곧 밀도 $x^{a-1}(1-x)^{b-1}/B(a,b)$($0<x<1$)를 따르고 $X+Y$와 독립이다.
\end{proposition}

\begin{proof}
$s=1$이라 해도 된다. 사상 $(x,y)\mapsto(u,v)=(x+y,\,x/(x+y))$의 역은 $x=uv$, $y=u(1-v)$이고 야코비안 행렬식의 절댓값은 $u$이다. 결합 밀도는
\[
\frac{(uv)^{a-1}(u(1-v))^{b-1}e^{-u}}{\Gamma(a)\Gamma(b)}\,u=\frac{u^{a+b-1}e^{-u}}{\Gamma(a+b)}\cdot\frac{v^{a-1}(1-v)^{b-1}}{B(a,b)}
\]
로 $u$의 함수와 $v$의 함수의 곱이다. 그러므로 두 변수는 독립이고 각각 $\mathrm{Gamma}(a+b)$와 $\mathrm{Beta}(a,b)$를 따른다.
\end{proof}

이 명제로 $F$와 베타가 이어진다. $F=(U/m)/(V/n)$이면 $mF/(mF+n)=U/(U+V)$이고 $U$와 $V$는 척도 $2$인 감마이므로 이것은 $\mathrm{Beta}(m/2,n/2)$이다. 따라서
\[
P(F\le x)=I_{mx/(mx+n)}\Bigl(\frac m2,\frac n2\Bigr)
\]
이고 같은 방법으로 $t$의 누적분포함수도 불완전 베타 함수로 쓰인다. 감마 확률변수의 역수의 분포인 \term{역감마}{inverse gamma}는 정규 모형의 분산에 대한 켤레 사전분포로 쓰이며 $X=\text{scale}/W$($W\sim\mathrm{Gamma}(a,1)$)로 정의된다.

\section{이산 분포}
\label{sec:discrete-families}

\subsection{이항, Poisson, 그리고 두 항등식}

성공 확률 $p$인 독립 시행 $n$번의 성공 횟수가 \term{이항 분포}{binomial distribution} $\mathrm{Bin}(n,p)$이고 $n=1$이 Bernoulli 분포이다. $n\to\infty$, $p=\mu/n$이면
\[
\binom nk\Bigl(\frac\mu n\Bigr)^k\Bigl(1-\frac\mu n\Bigr)^{n-k}=\frac{\mu^k}{k!}\cdot\frac{n(n-1)\cdots(n-k+1)}{n^k}\Bigl(1-\frac\mu n\Bigr)^{n-k}\longrightarrow\frac{\mu^ke^{-\mu}}{k!}
\]
로 \term{Poisson 분포}{Poisson distribution}가 나온다. 드문 사건의 횟수를 세는 분포이다. 근사의 정확도는 Le Cam의 부등식\cite{lecam1960}이 총변동 거리 $\sup_A|P(A)-Q(A)|$, 곧 $\sum_k|p_k-q_k|$의 절반으로 $np^2=\mu^2/n$ 이하라고 보장한다. $\mu=3$에서 실제 총변동 거리는 $n=10,100,1000,10^4$에서 $8.6\times10^{-2}$, $7.6\times10^{-3}$, $7.5\times10^{-4}$, $7.5\times10^{-5}$로 $1/n$에 비례하고 $n\ge100$에서 경계 $\mu^2/n$의 $1/12$쯤이다.

두 분포의 누적분포함수는 40장의 불완전 감마·베타 함수와 정확히 같다.

\begin{proposition}
\label{prop:cdf-identities}
정수 $k\ge0$에 대해
\[
P\bigl(\mathrm{Poisson}(\mu)\le k\bigr)=Q(k+1,\mu)
\]
이고, 정수 $0\le k\le n-1$에 대해
\[
P\bigl(\mathrm{Bin}(n,p)\le k\bigr)=I_{1-p}(n-k,\,k+1)
\]
이다($k\ge n$이면 이항의 확률은 $1$이다). $Q$는 정규화된 위쪽 불완전 감마 함수, $I$는 정규화된 불완전 베타 함수이다.
\end{proposition}

\begin{proof}
첫째 식. $Q(k+1,\mu)=\Gamma(k+1)^{-1}\int_\mu^\infty t^ke^{-t}\,dt$를 부분 적분하면 $Q(k+1,\mu)=\mu^ke^{-\mu}/k!+Q(k,\mu)$이고 $Q(1,\mu)=e^{-\mu}$이므로 귀납법으로 $Q(k+1,\mu)=\sum_{j\le k}\mu^je^{-\mu}/j!$이다.

둘째 식은 확률적으로 보는 것이 간명하다. $(0,1)$의 독립 균등 확률변수 $n$개 가운데 $p$ 이하인 것의 개수가 $\mathrm{Bin}(n,p)$이다. 개수가 $k$ 이하라는 것은 $(k+1)$번째로 작은 값이 $p$보다 크다는 것과 같다. 균등 표본의 $j$번째 순서통계량은 $\mathrm{Beta}(j,n-j+1)$을 따르므로(57장), 이 확률은 $1-I_p(k+1,n-k)=I_{1-p}(n-k,k+1)$이다. 마지막 등호는 $\mathrm{Beta}(a,b)$와 $1-\mathrm{Beta}(b,a)$가 같은 분포라는 데서 나온다.
\end{proof}

이 항등식이 통계 라이브러리가 이산 분포를 계산하는 방법이다. 이항의 누적분포함수를 $k+1$개 항의 합으로 계산하면 $n$이 클 때 느리고 위쪽 꼬리를 $1$에서 빼서 구하면 소거로 자릿수를 잃는다. 불완전 베타 함수로 계산하면 40장의 영역 분할(평균 근처에서는 균등 점근 전개)이 $n$이 커도 크게 늘지 않는 비용으로 양쪽 꼬리의 상대 정확도를 준다. 반대로 감마 분포에서 모양이 정수인 경우의 꼬리는 유한합으로 계산할 수 있어 stardust는 그것을 따로 쓴다.

\subsection{대기 시간과 혼합}

첫 성공까지의 시행 횟수가 \term{기하 분포}{geometric distribution}이고($P(X=k)=(1-p)^{k-1}p$, $k\ge1$) $r$번째 성공 전의 실패 횟수가 \term{음이항 분포}{negative binomial distribution}
\[
P(X=k)=\frac{\Gamma(k+r)}{k!\,\Gamma(r)}p^r(1-p)^k,\qquad k\ge0
\]
이다. 이 식은 정수가 아닌 $r>0$에도 뜻이 있는데, 그때의 해석은 혼합이다. $\Lambda\sim\mathrm{Gamma}(r,(1-p)/p)$를 뽑고 $X\mid\Lambda\sim\mathrm{Poisson}(\Lambda)$이면
\[
P(X=k)=\int_0^\infty\frac{\lambda^ke^{-\lambda}}{k!}\cdot\frac{\lambda^{r-1}e^{-\lambda p/(1-p)}}{\Gamma(r)((1-p)/p)^r}\,d\lambda=\frac{\Gamma(k+r)}{k!\,\Gamma(r)}p^r(1-p)^k
\]
이다(적분은 감마 적분이다). 그래서 음이항은 평균이 개체마다 다른 Poisson이고 분산이 평균보다 큰 과산포 계수 자료의 모형으로 쓰인다(63장). 같은 식으로 성공 확률을 베타에서 뽑은 이항이 \term{베타이항 분포}{beta-binomial distribution}이고 $t$ 분포는 분산을 역감마에서 뽑은 정규이다.

그 밖에 stardust는 비복원 추출의 성공 횟수인 \term{초기하 분포}{hypergeometric distribution}(58장의 Fisher 정확 검정의 귀무 분포), 정수 구간 $[\text{low},\text{high})$의 \term{이산 균등 분포}{discrete uniform distribution}, $P(X=k)=-p^k/(k\log(1-p))$인 로그 급수 분포, $P(X=k)=k^{-a}/\zeta(a)$인 Zipf 분포를 둔다. 로그 급수는 음이항에서 $0$을 빼고 $r\to0$으로 보낸 극한이며(음이항의 $1-p$가 로그 급수의 $p$가 된다) Zipf는 순위와 빈도의 멱법칙을 나타낸다.

\section{극값과 수명}
\label{sec:extreme-value}

독립이고 같은 분포를 따르는 $X_1,\dots,X_n$의 최댓값 $M_n$의 분포는 $F^n$이다. 적당한 상수 $a_n>0$, $b_n$으로 $(M_n-b_n)/a_n$이 퇴화하지 않는 분포로 수렴한다면, 그 극한은 위치와 척도를 빼면 다음 꼴 가운데 하나라는 것이 Fisher--Tippett--Gnedenko 정리이다\cite{fisher1928extreme,gnedenko1943}.
\[
G(z)=\exp\bigl(-(1+\xi z)^{-1/\xi}\bigr),\qquad1+\xi z>0
\]
$\xi=0$은 극한 $\exp(-e^{-z})$로 읽는다. 이것을 \term{일반화 극값 분포}{generalized extreme value distribution}라 하고 $\xi=0$이 Gumbel(정규·지수·로그정규처럼 꼬리가 어떤 멱보다도 빨리 줄 때), $\xi>0$이 Fréchet(꼬리가 멱으로 줄 때), $\xi<0$이 위로 유계인 역Weibull이다. 지수 분포의 최댓값에서 $b_n=\log n$, $a_n=1$로 두면 $P(M_n-\log n\le z)=(1-e^{-z}/n)^n\to\exp(-e^{-z})$로 Gumbel이 바로 보인다. stardust의 \code{GEV}는 scipy를 따라 모양을 $c=-\xi$로 쓰므로 $c>0$이 위로 유계인 쪽이다. 홍수의 연 최대 수위나 금융 손실의 꼬리가 이 분포로 모형화된다.

\term{베이불 분포}{Weibull distribution} $P(X>x)=\exp(-(x/s)^c)$는 최솟값의 극한 분포이자 수명 분포의 표준이다. 위험률 $f/(1-F)=(c/s)(x/s)^{c-1}$이 $c<1$이면 줄고 $c>1$이면 늘어 초기 고장과 마모를 모두 나타낸다. $c=1$이 지수, $c=2$가 레일리 분포이다(척도 $s$인 베이불은 척도 $s/\sqrt2$인 레일리이다). \term{파레토 분포}{Pareto distribution} $P(X>x)=(s/x)^b$는 소득이나 도시 크기처럼 꼬리가 멱으로 주는 양을 나타내며 $b$차 이상의 모멘트가 없다. 로지스틱 분포는 누적분포함수가 로지스틱 함수 $1/(1+e^{-z})$이어서 로지스틱 회귀(63장)의 잠재 변수 해석에 나타나고 라플라스 분포는 밀도가 $e^{-|z|}/2$로 중앙값 회귀(66장)의 우도이다. 코시 분포는 두 독립 표준 정규의 비이고 평균이 없으며 표본평균이 하나의 관측과 같은 분포를 따르므로 큰 수의 법칙이 성립하지 않는 반례이다(51장).

\section{비중심 분포}
\label{sec:noncentral}

$Z_i\sim N(\mu_i,1)$이 독립이면 $\sum Z_i^2$은 \term{비중심 카이제곱 분포}{noncentral chi-square distribution}를 따르며 그 분포는 자유도 $k$와 비중심도 $\lambda=\sum\mu_i^2$에만 의존한다(직교 변환으로 평균 벡터를 첫 축에 놓을 수 있다). 검정의 대립가설 아래에서 검정 통계량이 정확히 또는 점근적으로 따르는 분포이므로 검정력과 표본 크기 계산(55장)에 쓰인다.

\begin{proposition}
\label{prop:noncentral-mixture}
$\chi'^2_k(\lambda)$는 Poisson 혼합 $\sum_{j\ge0}e^{-\lambda/2}(\lambda/2)^j/j!\cdot\chi^2_{k+2j}$이다. 곧 $J\sim\mathrm{Poisson}(\lambda/2)$를 뽑고 $\chi^2_{k+2J}$를 뽑은 것과 같은 분포이다.
\end{proposition}

\begin{proof}
적률생성함수로 보인다(51장에서 적률생성함수가 분포를 결정함을 본다). $Z\sim N(\mu,1)$이면 $t<1/2$에서
\[
E\,e^{tZ^2}=\int\frac{e^{tz^2-(z-\mu)^2/2}}{\sqrt{2\pi}}\,dz=(1-2t)^{-1/2}\exp\Bigl(\frac{\mu^2t}{1-2t}\Bigr)
\]
이므로(지수의 완전제곱), 합의 적률생성함수는 $(1-2t)^{-k/2}\exp(\lambda t/(1-2t))$이다. 한편 혼합의 적률생성함수는
\[
\sum_je^{-\lambda/2}\frac{(\lambda/2)^j}{j!}(1-2t)^{-(k+2j)/2}=(1-2t)^{-k/2}\exp\Bigl(-\frac\lambda2+\frac{\lambda/2}{1-2t}\Bigr)
\]
이고 지수가 $\lambda t/(1-2t)$로 같다.
\end{proof}

그래서 비중심 카이제곱의 누적분포함수는 Poisson 가중치를 붙인 불완전 감마 함수의 급수이다. 가중치는 $j\approx\lambda/2$ 둘레에 몰려 있으므로 합은 $j=0$이 아니라 가장 큰 가중치에서 시작해 양쪽으로 더해 나가는 것이 좋다. 자유도 $5$, $\lambda=100$, $x=120$에서 누적분포함수 $0.77731$을 상대 $10^{-17}$ 수준까지 더하려면 $j=4$부터 $101$까지 $98$개 항이 필요하다. 가중치의 표준편차가 $\sqrt{50}\approx7$이어서 아래쪽은 그 여섯 배 남짓 바깥까지가 유효하고, 위쪽은 곱해지는 $P(\chi^2_{5+2j}\le x)$가 빨리 줄어 더 일찍 끝난다. 비중심 $F$도 같은 Poisson 혼합으로 불완전 베타 함수의 급수가 되고, 비중심 $t$는 Poisson형 가중치가 붙은 불완전 베타 함수의 두 급수가 된다.

\section{서른한 개의 목록}
\label{sec:family-list}

\cref{tab:continuous-families}와 \cref{tab:discrete-families}가 stardust의 단변량 분포이다. 매개변수의 이름은 stardust의 것이다.

\begin{table}[t]
\begin{sdtable}{colspec={llll}, row{1}={font=\small\bfseries}, cells={font=\small}}
\toprule
이름 & 매개변수 & 지지집합 & 정의 \\
\midrule
\code{Normal} & \code{mu}, \code{sigma} & $\mathbb R$ & $\phi(z)/\sigma$ \\
\code{StudentT} & \code{df} & $\mathbb R$ & $Z/\sqrt{V/\nu}$ \\
\code{ChiSquared} & \code{df} & $(0,\infty)$ & $\mathrm{Gamma}(k/2,2)$ \\
\code{F} & \code{dfn}, \code{dfd} & $(0,\infty)$ & $(U/m)/(V/n)$ \\
\code{Gamma} & \code{a}, \code{scale} & $(0,\infty)$ & $x^{a-1}e^{-x/s}/(\Gamma(a)s^a)$ \\
\code{Beta} & \code{a}, \code{b} & $[0,1]$ & $x^{a-1}(1-x)^{b-1}/B(a,b)$ \\
\code{Exponential} & \code{scale} & $[0,\infty)$ & $e^{-x/s}/s$ \\
\code{LogNormal} & \code{mu}, \code{sigma} & $(0,\infty)$ & $\log X\sim N(\mu,\sigma^2)$ \\
\code{Uniform} & \code{a}, \code{b} & $[a,b]$ & $1/(b-a)$ \\
\code{Cauchy} & \code{mu}, \code{sigma} & $\mathbb R$ & $1/(\pi\sigma(1+z^2))$ \\
\code{Laplace} & \code{mu}, \code{sigma} & $\mathbb R$ & $e^{-|z|}/(2\sigma)$ \\
\code{Weibull} & \code{c}, \code{scale} & $[0,\infty)$ & $P(X>x)=e^{-(x/s)^c}$ \\
\code{Logistic} & \code{mu}, \code{sigma} & $\mathbb R$ & $F=1/(1+e^{-z})$ \\
\code{Gumbel} & \code{mu}, \code{sigma} & $\mathbb R$ & $F=\exp(-e^{-z})$ \\
\code{GEV} & \code{c}, \code{mu}, \code{sigma} & $1-cz>0$ & $F=\exp(-(1-cz)^{1/c})$ \\
\code{Pareto} & \code{b}, \code{scale} & $[s,\infty)$ & $P(X>x)=(s/x)^b$ \\
\code{Rayleigh} & \code{scale} & $[0,\infty)$ & $P(X>x)=e^{-x^2/(2s^2)}$ \\
\code{InverseGamma} & \code{a}, \code{scale} & $(0,\infty)$ & $s/\mathrm{Gamma}(a,1)$ \\
\code{NoncentralChiSquared} & \code{df}, \code{nc} & $(0,\infty)$ & $\sum(Z_i+\mu_i)^2$ \\
\code{NoncentralT} & \code{df}, \code{nc} & $\mathbb R$ & $(Z+\delta)/\sqrt{V/\nu}$ \\
\code{NoncentralF} & \code{dfn}, \code{dfd}, \code{nc} & $(0,\infty)$ & $(U'/m)/(V/n)$ \\
\bottomrule
\end{sdtable}
\caption{stardust의 단변량 연속 분포 스물한 개. $z=(x-\mu)/\sigma$이고 $s$는 \code{scale}이다. $U'$은 비중심 카이제곱이다. 비중심 분포의 $Z$, $Z_i$는 표준 정규이고 \code{nc}는 $\lambda=\sum\mu_i^2$ 또는 $\delta$이다.}
\label{tab:continuous-families}
\end{table}

\begin{table}[t]
\begin{sdtable}{colspec={llll}, row{1}={font=\small\bfseries}, cells={font=\small}}
\toprule
이름 & 매개변수 & 지지집합 & 확률 질량 함수 \\
\midrule
\code{Bernoulli} & \code{p} & $\{0,1\}$ & $p^k(1-p)^{1-k}$ \\
\code{Binomial} & \code{n}, \code{p} & $\{0,\dots,n\}$ & $\binom nkp^k(1-p)^{n-k}$ \\
\code{Poisson} & \code{mu} & $\{0,1,\dots\}$ & $\mu^ke^{-\mu}/k!$ \\
\code{Geometric} & \code{p} & $\{1,2,\dots\}$ & $(1-p)^{k-1}p$ \\
\code{NegativeBinomial} & \code{n}, \code{p} & $\{0,1,\dots\}$ & $\Gamma(k+n)p^n(1-p)^k/(k!\Gamma(n))$ \\
\code{Hypergeometric} & \code{total}, \code{successes}, \code{draws} & 정수 구간 & $\binom Kk\binom{N-K}{d-k}/\binom Nd$ \\
\code{DiscreteUniform} & \code{low}, \code{high} & $\{\text{low},\dots,\text{high}-1\}$ & $1/(\text{high}-\text{low})$ \\
\code{BetaBinomial} & \code{n}, \code{a}, \code{b} & $\{0,\dots,n\}$ & $\binom nkB(k+a,n-k+b)/B(a,b)$ \\
\code{LogSeries} & \code{p} & $\{1,2,\dots\}$ & $-p^k/(k\log(1-p))$ \\
\code{Zipf} & \code{a} & $\{1,2,\dots\}$ & $k^{-a}/\zeta(a)$ \\
\bottomrule
\end{sdtable}
\caption{stardust의 단변량 이산 분포 열 개. 초기하에서 $N$은 \code{total}, $K$는 \code{successes}, $d$는 \code{draws}이고 지지집합은 $\max(0,d-N+K)$부터 $\min(K,d)$까지이다.}
\label{tab:discrete-families}
\end{table}

\section{분포 계산의 수치 기법}
\label{sec:families-numerics}

분포를 특수 함수의 조합으로 쓰면 조합하는 자리에서 자릿수를 잃을 수 있다. 정규 밀도 $\exp(-z^2/2)$에서 $z=(x-\mu)/\sigma$를 binary64로 반올림하면 상대오차 $u$가 지수에서 $z^2$배로 커져 $z=30$이면 결과의 상대오차가 $900u$이다(39장에서 본 $\exp(-x^2)$의 문제와 같다). 감마의 $x/s$, 베이불의 $c\log(x/s)$, 로그정규의 $\log x$도 같다. 또 $t$와 $F$의 누적분포함수는 불완전 베타 함수의 인자 $\nu/(\nu+t^2)$이 $1$에 가까워지는 곳에서 $1-x$를 직접 계산해야 하고 비중심 분포는 앞 절처럼 가중치의 봉우리에서 급수를 시작해야 한다. 정확한 분포 계산은 결국 40장까지의 특수 함수를 언제, 어떤 인자로 부를지의 문제이다.

\begin{instardust}[\code{docs/design/stats.md} §1.2, §3.1, T25, T26, T31]
stardust의 분포는 연속 21개(기본 12, 극값·수명 6, 비중심 3)와 이산 10개이며 다변량 7개는 52장에서 다룬다. 몇 가지 정의는 문헌마다 다르므로 설계 문서가 정했다. \code{Geometric}은 첫 성공까지의 시행 횟수($k\ge1$), \code{NegativeBinomial}은 $n$번째 성공 전의 실패 횟수($k\ge0$), \code{DiscreteUniform}(\code{low}, \code{high})는 반열린 구간이고(난수의 \code{integers}와 같다) \code{GEV}의 모양은 scipy처럼 $c=-\xi$이다. 자유도가 무한대이면 극한 분포로 정의한다. \code{StudentT(inf)}는 \code{Normal()}, \code{F(dfn, inf)}는 $\chi^2_{\text{dfn}}/\text{dfn}$이고 밀도부터 모멘트와 표본까지 모든 메서드가 그 극한의 값을 낸다. 표준화 변수는 정확한 입력에서 double-double로 만들고 지수 $-z^2/2-\log(\sigma\sqrt{2\pi})$를 double-double로 만들어 한 번에 지수화한다. 누적분포함수는 반올림한 $z$에서 계산한 뒤 $f(z)\cdot z_{\mathrm{lo}}$로 1차 보정한다.
\end{instardust}

\section{맺음}

분포족은 몇 가지 구성으로 이어진다. 정규의 제곱합과 비에서 카이제곱, $t$, $F$가 나오고 감마의 합과 비에서 베타가 나오며 극한에서 Poisson과 극값 분포가, 혼합에서 음이항과 베타이항이 나온다. 이 관계는 이론의 정리일 뿐 아니라 계산의 처방이다. 이항과 Poisson, $t$와 $F$, 카이제곱의 누적분포함수가 모두 불완전 베타·감마 함수로 바뀌므로, 통계 라이브러리의 정확도는 40장의 두 함수의 정확도로 모인다. 다음 장에서는 분포를 모멘트와 생성함수로 요약하고 큰 수의 법칙과 중심극한정리로 분포의 극한을 다룬다.

분포의 성질과 관계는 Johnson, Kotz, Balakrishnan\cite{johnson1994continuous,johnson1995continuous}과 Johnson, Kemp, Kotz\cite{johnson2005discrete}의 백과사전적인 책들과 Casella와 Berger의 수리통계 교과서\cite{casella2002}에 자세하다.

\chapterreferences
